\section{Introduction}

\DraftingNote{Motivate environmental sound classification under real background
noise. Introduce the gap between clean test accuracy and controlled robustness,
then summarize the study contributions. Add only verified primary citations.}

This study investigates whether training-time audio augmentation reduces
performance degradation when urban sound classifiers are evaluated with
controlled background noise. Three architectures are compared: a compact
two-dimensional convolutional neural network (CNN), a convolutional recurrent
neural network (CRNN) with a bidirectional gated recurrent unit (BiGRU), and a
non-pretrained ResNet18 adapted to one-channel Log-Mel spectrograms.

\subsection{Research Questions}

\begin{enumerate}
  \item Does training-time audio augmentation reduce the loss of classification
  performance when test audio is corrupted by background noise?
  \item Is the augmentation effect consistent across CNN, CRNN, and ResNet18
  architectures?
\end{enumerate}

\subsection{Contributions}

\DraftingNote{Refine this list after the rest of the paper is stable. Keep the
claims within the one-fold, one-training-seed experimental boundary.}

\begin{itemize}
  \item A reproducible audio-classification pipeline with fold-aware data
  handling, shared preprocessing, and common model interfaces.
  \item A controlled comparison of baseline and augmented training across three
  architectures and four held-out test conditions.
  \item Deterministic noise corruption with separated training and evaluation
  noise banks and per-sample provenance.
  \item Architecture-level, robustness-curve, and per-class analysis of the
  measured augmentation effects.
\end{itemize}
